%!TEX program = uplatex
\RequirePackage{plautopatch}
\documentclass[uplatex,dvipdfmx]{jlreq}
\usepackage{amsmath,amssymb,amscd,amsthm}

\newtheorem{thm}{定理}
\newtheorem{theorem}[thm]{定理}
\newtheorem{prop}[thm]{命題}
\newtheorem{proposition}[thm]{命題}
\newtheorem{cor}[thm]{系}
\newtheorem{corollary}[thm]{系}
\newtheorem{lem}[thm]{補題}
\newtheorem{lemma}[thm]{補題}
\newtheorem{conj}[thm]{Conjecture}
\newtheorem{remark}[thm]{注意}
\newtheorem{defi}[thm]{定義}
\newtheorem{definition}[thm]{定義}
\usepackage{url} % or hyperref
\pagestyle{empty}

\begin{document}

%======================================================================
% 講演１：楠岡成雄「ブラウン運動から確率微分方程式まで」
%======================================================================
\begin{center}
  {\Large\bfseries ブラウン運動から確率微分方程式まで}
\end{center}

\bigskip
\hspace*{\fill}{\large\bfseries 楠岡　成雄（東大・数理）}

\bigskip

講演では、まず離散時間の確率過程モデルについて述べ、
それを連続時間モデルにするにあたってどのような困難があったかを述べる。
ついで、ブラウン運動の定義と構成について詳しく述べる。
ついで、確率微分方程式とは何かについて、
歴史的ないきさつを交えて述べる。
確率積分に対する Bacherier, L\'evy らの直感的なアイデア、
Kolmogorov の拡散方程式の考え方について述べ、
確率積分の定義等について述べる。
最後に時間があれば
Malliavin 解析などの
確率解析の全体像やその応用の可能性について述べたい。

講演では、確率論の基礎知識を仮定する。
数学的に厳密な証明を述べることはほとんどできないが、
確率解析をこれから勉強しようとする人のための
道案内となるような講演をしたいと考えている。

以下では、日本語の参考文献をいくつかあげておく。

\begin{thebibliography}{99}
\bibitem{funaki}
舟木直久，確率微分方程式，岩波講座現代数学の基礎，岩波書店，
1997
\bibitem{nagai}
長井英生，確率微分方程式，共立出版，1999
\bibitem{Watanabe1975}
渡辺信三，確率微分方程式，産業図書，1975
\end{thebibliography}

\newpage

%======================================================================
% 講演２：重川一郎「マリアバン解析について」
%======================================================================
\begin{center}
  {\Large\bfseries マリアバン解析について\\
  \normalsize --- Brown 運動の汎関数を解析する}
\end{center}

\bigskip
\hspace*{\fill}{\large\bfseries 重川　一郎（京都大学大学院理学研究科）}

\bigskip

世の中の現象を眺め渡してみると，
日が昇ったり，沈んだりの規則的な現象も多いが，
株が上がったり，下がったりの偶然的な現象も結構あることに気づく．
量子力学はそもそも確率的にしか予測できない，とかいう話もあって，
偶然現象というのはむしろ規則的なものより根が深いのかもしれない．
そういう偶然的な現象を数学的に記述しようとするとき，
Brown 運動がよく用いられる．
つまり現象が Brown 運動の汎関数として表現されるということである．
ここでは，例えば確率微分方程式の解などというものを
典型的なものとして考えている訳である．

そこで次のことを問題にする：
\begin{center}
Brown 運動の汎関数を解析するにはどうしたらよいのか？
\end{center}
この答えが一つであるはずはないが，一つの答え方として，
\begin{center}
Brown 運動に対する{\bfseries 微積分}を展開する
\end{center}
ということが可能である．
実際それを実行したのが Paul Malliavin で，
それは1970年代の後半のことであった．
この理論が出た当初は，センセーショナルな事件であったが，
今では常識とされるほど基本的なものとなっている．

Brown 運動は関数空間の上の測度として実現される．
測度があると，「積分」が自然に定義できる．
また Brown 運動が実現されている関数空間は Banach 空間なので，
無限次元ではあっても Fr\'echet 微分というものがすでに定義されている．
しかしこの微分概念は確率微分方程式の解のようなものを考えると
有効なものではないことが分かる．
実際確率微分方程式の解は， Fr\'echet の意味で微分できないばかりでなく，
連続ですらない．
この困難のために，微分の概念を拡張することが求められたのであったが，
中でも次のようなことは満たされていることが望ましい．
\begin{enumerate}
\item 確率微分方程式の解が微分できる
\item Brown 運動の積分と調和している
\item 実際に計算可能である
\end{enumerate}
特に，2番目の条件が重要で，これなくしては微積分とは言い難い．
微分と積分がバラバラの方向を向いていたのでは実り豊かな結果は得られない．
この微分と積分が調和しているというのは，実際には部分積分の公式として
表現できる．
また，微分とは変分を考えることであるが，その変分の取り方が測度と
密接に関連している必要がある．
そのためにずらしに対する絶対連続性が
重要な位置を占め，いわゆる Cameron-Martin 空間と呼ばれる空間が
鍵となってくる．

Brown 運動の測度は Gauss 測度と呼ばれる範疇に入る．
講演では Gauss 測度の話を進めながら，それと微分をどう調和させていくか，
という方向性で話を進めていく．

\begin{thebibliography}{9}
\bibitem{Kusuoka84} 楠岡 成雄,
Malliavin calculus とその応用,
{\it 数学}, {\bfseries 36} (1984), 97--109.
\bibitem{Kusuoka93} 楠岡 成雄,
無限次元解析としての確率解析,
{\it 数学}, {\bfseries 45} (1993), 289--298.
\bibitem{Nualart95} D.~Nualart,
``{\it The Malliavin calculus and related topics\/},''
Probability and its Applications, Springer-Verlag, New York, 1995.
\bibitem{Shigekawa98} 重川 一郎,
確率解析, 「岩波講座現代数学の展開」, 岩波書店, 東京, 1998.
\bibitem{Watanabe84} S.~Watanabe,
``{\it Lectures on stochastic differential equations and Malliavin
calculus\/},'' Tata Institute of Fundamental Research,
Springer-Verlag, Berlin-Heidelberg-New York, 1984.
\bibitem{Watanabe90} 渡辺 信三,
確率解析とその応用,
{\it 数学}, {\bfseries 45} (1990), 97--110.
\end{thebibliography}

\newpage

%======================================================================
% 講演３：谷口説男「確率解析と幾何」
%======================================================================
\begin{center}
  {\Large\bfseries 確率解析と幾何\\
  \normalsize --- ユビキタス・ウィナー・インテグラル ---}
\end{center}

\bigskip
\hspace*{\fill}{\large\bfseries 谷口　説男（九州大学大学院数理学研究院）}

\bigskip

この講演では確率解析が幾何学に絡んでくるいくつかの局面を紹介
いたします．

\begin{center}
\begin{minipage}{11cm}
『ありとあらゆるものが何らかの確率変数，すなわち予測不可能な
量を記述する関数の積分として表示することができる』
\end{minipage}
\end{center}
というのがモットーです．
些か奇異な感じがするかも知れませんが，確率解析が解析学や幾何
学に顔を出すときにはあながち的外れなことではありません．
この講演で触れるいくつかの量は，リーマン多様体上の経路を変数
とする関数を経路空間(経路全体のなす空間)上での積分することで
得られます．
測地座標により多様体とユークリッド空間が対応していることに呼
応して，リーマン多様体に値をとる経路を変数とする関数の積分は，
ユークリッド空間に値をとる経路を変数とする関数の積分へと置き
換えられます．
標題にあるウィナー・インテグラルとは，このようなユークリッド
空間値経路を変数とする関数の積分のことです．

経路変数の関数の関数の積分といえば，ファインマン経路積分が思
い起こされます．
ファインマン経路積分論は，シュレディンガー作用素に附随するプ
ロパゲータは，作用積分を指数とする指数関数の経路全般に渡る積
分(ファインマン経路積分)として表示されるといいます．
準古典近似により古典力学が再現される様が経路空間上の停留位相
法として了解されたりと，この表示の示唆するところは非常に奥深
いものがあります．
残念なことにこの積分は仮想的なもので数学的厳密さは欠けていま
すが，その時間分割近似による定義は，作用積分の係数が純虚数な
ことを除いてウィナー積分の定義と全く同じ形をしています．
ウィナー積分は数学的厳密さを持って定義されていますから，この
類似性からウィナー積分に関わる種々の漸近理論の考察など興味深
い問題が生起してきます．
このような方向のウィナー積分研究は，M.~Kacが1947年にコーネル
大学物理学講演会でR.~Feynmanの博士論文に関する発表を聞いたと
きに始まり，その後大偏差原理，確率論的漸近展開理論などへと発
展しました．

リーマン多様体を調べるひとつの切口として，ラプラシアン，もし
くは有界作用素ではないラプラシアンよりも扱い易いラプラシアン
に附随する半群(熱半群)があります．
上で触れたファインマン経路積分との酷似から予想されるように，
熱半群はウィナー積分として表すことができ，これによりリーマン
多様体を確率論を用いて解析する手がかりを得ることができます．
このときリーマン多様体値の経路を変数とする関数を積分しますが，
どのような測度に関して積分するのかはまだ明らかでありません．
マリアバン解析についての講演の中で明らかになるように，ユーク
リッド空間で確率解析を展開するときはガウス測度とガウス測度か
ら作られるブラウン運動が基本的な役割を果たします．
ブラウン運動はユークリッド空間の基本的なランダムな曲線を記述
しているともいえます．
ユークリッド空間の基本的な曲線は直線ですが，リーマン多様体上
のそれは測地線です．
ユークリッド空間の直線がリーマン多様体に巻き付けられて測地線
となるように，ユークリッド空間上のブラウン運動を巻き付けるこ
とでリーマン多様体上の基本的なランダムな曲線であるブラウン運
動が得られます．
このリーマン多様体上のブラウン運動が導く測度について積分を行
うことで熱核がウィナー積分として表示され，この表示から，熱核
の短時間漸近挙動，指数定理の熱方程式的証明の確率解析による解
釈などが可能となってきます．
またブラウン運動をランダムな測地線ととらえればその写像として
の性質に曲率の反映を見いだすことができます．

この講演では
\begin{enumerate}
\item リーマン多様体上のブラウン運動の構成
\item 曲率とブラウン運動
\item 確率解析による熱核の漸近展開
\item Gauss-Bonnet-Chernの定理の確率解析による証明
\end{enumerate}
について紹介します．
時間があれば，幾何を離れますが，KdV方程式の$n$ソリトン解がウィ
ナー積分として表示されることを紹介したいと思います．

講演用の詳しいノートを
\begin{center}
\url{http://www.math.kyushu-u.ac.jp/~taniguch/}
\end{center}
の``Preprints and so on''で公開する予定です(12月初旬)．
興味をお持ちの方はご覧下さい．
教科書としては以下のものを挙げておきます．
詳しい文献表は上のノートに付けます．

\begin{thebibliography}{99}
\bibitem{kde1}
K. D. Elworthy,
Stochastic differential equations on manifolds,
Cambridge Univ. Press, 1982.
\bibitem{kde2}
K. D. Elworthy,
Geometric aspects of diffusions on manifolds,
Ecole d'Et\'e Probabilit\'es de Saint-Flour XV-XVI 1987,
1989,
P.L. Henequin ed., Lect. Notes in Math. {\bfseries 1362}, 1989,
Springer, 276--425.
\bibitem{iw}
N. Ikeda and S. Watanabe,
Stochastic differential equations and diffusion processes,
2nd ed.,
North-Holland/Kodansha, Amsterdam/Tokyo, 1989.
\bibitem{dan}
D. W. Stroock,
An introduction to the analysis of paths on a Riemannian
manifold,
AMS, 2000.
\end{thebibliography}

\end{document}
